\chapter*{参考书目}
\addcontentsline{toc}{chapter}{参考书目}

\begin{enumerate}[leftmargin=2.6em,itemsep=0.55em]
\item G. B. Arfken, H. J. Weber, F. E. Harris, \emph{Mathematical Methods for Physicists}, 7th ed., Academic Press, 2013.
\item M. L. Boas, \emph{Mathematical Methods in the Physical Sciences}, 3rd ed., Wiley, 2005.
\item K. F. Riley, M. P. Hobson, S. J. Bence, \emph{Mathematical Methods for Physics and Engineering}, 3rd ed., Cambridge University Press, 2006.
\item E. Butkov, \emph{Mathematical Physics}, Addison-Wesley, 1968.
\item R. Courant, D. Hilbert, \emph{Methods of Mathematical Physics}, Vol. I--II, Wiley, 1953.
\item P. M. Morse, H. Feshbach, \emph{Methods of Theoretical Physics}, McGraw-Hill, 1953.
\item J. D. Jackson, \emph{Classical Electrodynamics}, 3rd ed., Wiley, 1999.
\item L. D. Landau, E. M. Lifshitz, \emph{Mechanics}, 3rd ed., Butterworth-Heinemann, 1976.
\item G. N. Watson, \emph{A Treatise on the Theory of Bessel Functions}, 2nd ed., Cambridge University Press, 1944.
\item M. Abramowitz, I. A. Stegun, \emph{Handbook of Mathematical Functions}, Dover, 1965.
\item 顾樵，《数学物理方法》，科学出版社，2012.
\item 梁昆淼，《数学物理方法》，第五版，高等教育出版社，2010.
\end{enumerate}

\bigskip
\noindent\small 说明：上述教材在选题传统与讲法上给了本书很多启发，尤其是顾樵对物理动机的强调与 Arfken 对方法适用条件的提醒。本书没有整段翻译或照搬其中任何内容；所有例题、数据与习题均为独立编写。
